\documentclass[10pt]{article}
\usepackage[T1]{fontenc}
\usepackage{mathpazo}
\usepackage[a5paper,margin=16mm]{geometry}
\usepackage{amsmath,amsthm,mathtools}
\usepackage{microtype}
\newtheorem{theorem}{Theorem}
\newtheorem{lemma}[theorem]{Lemma}
\theoremstyle{definition}
\newtheorem{definition}{Definition}
\DeclarePairedDelimiter{\ceil}{\lceil}{\rceil}
\DeclarePairedDelimiterX{\inner}[2]{\langle}{\rangle}{#1,\,#2}
\mathtoolsset{mathic=true}

\begin{document}
\section{Bounds from Cauchy--Schwarz}
\begin{definition}
The \emph{energy} of $v\in\mathbb{R}^{n}$ is $\mathcal{E}(v)\coloneqq\sum_{i=1}^{n}v_i^{2}$.
\end{definition}
\begin{lemma}[Cauchy--Schwarz]\label{lem:cs}
For all $u,v\in\mathbb{R}^{n}$,
\begin{equation}
  \inner{u}{v}^{2}\le \mathcal{E}(u)\,\mathcal{E}(v).
\end{equation}
\end{lemma}
\begin{proof}
Expand $0\le\sum_{i,j}(u_iv_j-u_jv_i)^{2}$ and divide by two.
\end{proof}
\begin{theorem}\label{thm:mean}
For positive reals $a_1,\dots,a_n$ the arithmetic mean dominates the geometric mean:
\begin{align}
  \frac{a_1+\dotsb+a_n}{n} &\ge \sqrt[n]{a_1\cdots a_n},\\
  \ceil*{\frac{n+1}{2}} &\le n \qquad (n\ge1).
\end{align}
\end{theorem}
\begin{proof}
Apply Lemma~\ref{lem:cs} with $u_i=\sqrt{a_i}$ and $v_i=1/\sqrt{a_i}$ to obtain
$n^{2}\le\bigl(\sum a_i\bigr)\bigl(\sum a_i^{-1}\bigr)$.
\end{proof}

\paragraph{Display variety.}
\begin{alignat*}{2}
  \int_0^{\pi}\sin x\,dx &= 2, &\qquad \sum_{k=1}^{\infty}\frac{1}{k^{2}} &= \frac{\pi^{2}}{6},\\
  \mathbb{E}[X] &= \smashoperator{\sum_{x\in\Omega}}x\,p(x), & \operatorname*{arg\,min}_{t} f(t) &\in \mathbb{R}.
\end{alignat*}
\[
  \begin{pmatrix} a & b\\ c & d\end{pmatrix}^{-1}
  =\frac{1}{ad-bc}\begin{pmatrix}\hphantom{-}d & -b\\ -c & \hphantom{-}a\end{pmatrix}
\]
\end{document}
